\documentclass[10pt,a4paper]{amsart}
\usepackage[margin=19mm]{geometry}
\usepackage{amsmath,amssymb,mathtools}
\usepackage[scr=rsfs,cal=euler]{mathalfa}
\usepackage{xfrac}
\usepackage[unicode,hidelinks,bookmarks=false]{hyperref}
\newcommand{\ie}{i.e.\ }
\newcommand{\cf}{cf.\ }
\newcommand{\resp}{respectively\ }
\newcommand{\Subsection}{\subsection}
\providecommand{\nobreakdash}{\nobreak\hbox{-}}
\setlength{\emergencystretch}{2em}
\multlinegap0pt
\usepackage{amssymb}
\usepackage{dsfont}
\newtheorem{theorem}{Theorem}
\newtheorem{lemma}{Lemma}
\newcommand{\C}{{\mathbb C}}
\newcommand{\N}{{\mathbb N}}
\newcommand{\R}{{\mathbb R}}
\newcommand{\Z}{{\mathbb Z}}
\newcommand{\abs}[1]{\left\vert#1\right\vert}
\DeclareMathOperator{\essinf}{ess\,inf}
\newcommand{\pa}[1]{\left(#1\right)}
\newcommand{\sinc}{\mathrm{sinc}}
\newcommand{\suc}[1]{\left\{#1\right\}}
\title{Dynamical sampling, derivatives, and interpolation formulas in the Paley--Wiener space}
\author{Iker Gardeazabal-Guti\'errez,, Mateus Sousa}
\date{Source: arXiv:2506.16195v2 (CC BY 4.0)}
\begin{document}
\maketitle

\noindent\emph{Source and licence.} Dynamical sampling, derivatives, and interpolation formulas in the Paley--Wiener space. Version arXiv:2506.16195v2 (CC BY 4.0), distributed by arXiv under the Creative Commons Attribution 4.0 International licence, http://creativecommons.org/licenses/by/4.0/, as recorded in the arXiv licence metadata for this exact version and shown on the abstract page https://arxiv.org/abs/2506.16195v2. Full licence notice: \texttt{licenses/arxiv-2506.16195-CC-BY-4.0.txt}.\par\smallskip

\noindent\textit{Editorial scope.} Counted unit: the article's theorem theo (source0.tex lines 283--309). The article's complete original proof of it is delivered with it (source0.tex lines 434--526, one \texttt{proof} environment of 93 lines, which the article places away from the statement). No mathematics is written, repaired or re-derived: the statement and the proof are the article's own lines, in the article's own notation, and no retained line is substituted or re-flowed.  Retained uncounted as the source-local context the proof needs: lines 367--380; lines 386--392.  Nothing was omitted from any retained span, so the package carries no omission record.  No retained line contains a citation, so the wrapper carries no bibliography and no bibliography entry.  No retained line contains a figure environment, an image-inclusion command or drawing code, so no image is delivered and the package declares no asset dependency.  Editorial formatting only, in addition: an AMS \texttt{amsart} preamble that restates the article's own declarations and macros used by the retained lines, namely the environments \texttt{theorem}, \texttt{lemma} on separate counters, together with the standard packages those lines need.  The retained environments print in this document as Theorem 1, Lemma 1 in order of appearance, whereas the article prints the same items with the numbering of its own full document.  The complete native source of the article as distributed in the arXiv e-print for this exact version is bundled unchanged as \texttt{sources/180134/source0.tex}.
\begin{lemma}\label{lem1}
	Let $M$ be a $N\times N$ matrix whose entries are $L^\infty(a,b)$ functions, and its associated linear operator
	\begin{align*}
		O_M: \pa{L^2(a,b)}^N &\to \pa{L^2(a,b)}^N \\
		(f_n)_{n=1}^N &\to M (f_n)_{n=1}^N.
	\end{align*}
	Then we have that: 
	
	- $O_M$ is injective if and only if $M(x)\neq 0$ for almost every $x\in (a,b)$,
	
	- $O_M$ is surjective if and only if $\essinf_{x\in (a,b)}\abs{\det M(x)}>0$,
	
	- $O_M$ is invertible if and only if $\essinf_{x\in (a,b)}\abs{\det M(x)}>0$.
\end{lemma}

\begin{lemma}\label{lem2}
    Let $T$ be a Fourier multiplier operator associated to a bounded and measurable multiplier $K$, $f$ a function 
    in $PW_\pi$, $N\in\N$ and $a\in\C$. Then we have that
	$$\sum_{m\in\Z} T(f)(Nm+a) e^{2\pi i m x}
	=\frac{1}{N} \sum_{m=1}^N \widehat{f}\pa{\frac{m-1-x}{N}} K\pa{\frac{m-1-x}{N}}e^\frac{2\pi i a (m-1-x)}{N},$$
	for $x\in \pa{\frac{N-2}{2},\frac{N}{2}}$.
\end{lemma}

\begin{theorem}\label{theo}
    Let $N\in\N$ and $T=\pa{T_n}_{n=1}^N$ where, for each $n=1,...,N$, $T_n$ is a Fourier multiplier operator associated 
    to a bounded  multiplier $K_n$. Let us define for $x\in\pa{\frac{N-2}{2},\frac{N}{2}}$ the $N \times N$ matrix
	$$\mathcal{M}_T(x)=\pa{\frac{1}{N} K_n\pa{\frac{m-1-x}{N}}}_{m,n=1}^N.$$
	Then the following are equivalent:
	
%	1) $\mathcal{M}_T\in L^\infty\pa{\pa{\frac{N-2}{2},\frac{N}{2}}, GL_N}$ and
%	$\pa{\mathcal{M}_T}^{-1}\in L^\infty\pa{\pa{\frac{N-2}{2},\frac{N}{2}}, GL_N}$,
%	
%	2) There exists $\epsilon>0$ such that $\abs{\det \mathcal{M}_T(x)} >\epsilon$ for almost every $x\in \pa{\frac{N-2}{2},\frac{N}{2}}$,
    1) $\essinf_{x\in \pa{\frac{N-2}{2},\frac{N}{2}}} \abs{\det \mathcal{M}_T(x)}>0$,
	
	2) There exists, for $n=1,...,k$, functions $g^T_n\in PW_\pi$ such that
	$$f(x)=\sum_{n=1}^N \sum_{m\in\Z} T_n(f)(Nm) g^T_n(x-Nm),$$
	for all $f\in PW_\pi$, where the convergence holds both in $L^2$ sense and uniformly.
	
	Furthermore, the functions $g^T_n$ appearing in 3) can be uniquely determined by either of the following 
	characterizations:
	
	i) For each $n=1,...,N$, the function $g^T_n$ is the unique function in $PW_\pi$ satisfying
	$$T_m(g^T_n)(Nj)=\delta_{n,m}\delta_{j,0},$$
	for all $m=1,...,N$ and $j\in\Z$,
	
	ii) The functions $g^T_n$, $n=1,...,N$, are the unique functions in $PW_\pi$ satisfying
	$$\pa{\widehat{g^T_n}\pa{\frac{j-1-\xi}{N}}}_{n,j=1}^N=\pa{\mathcal{M}_T(\xi)}^{-1},$$
	for almost every $\xi\in\pa{\frac{N-2}{2},\frac{N}{2}}$.
\end{theorem}

\begin{proof}[First proof of Theorem \ref{theo}:]
    We start by considering the operator $A_T: \pa{\ell^2}^N \to \pa{\ell^2}^N$ defined as
    $$A_T\pa{\pa{\left\{ b^n_m \right\}_{m\in\Z}}_{n=1}^N}
	=\pa{\left\{ \sum_{s=1}^N\sum_{j\in\Z} b_j^s T_n\pa{\sinc}(N(m-j)-s+1)\right\}_{m\in\Z}}_{n=1}^N,$$
	note that if $f\in PW_\pi$ we have that
	$$A_T\pa{\pa{\left\{ f(Nm+n-1) \right\}_{m\in\Z}}_{n=1}^N}
	=\pa{\left\{ T_n(f)(Nm)\right\}_{m\in\Z}}_{n=1}^N.$$
	We also consider the operator $F_N:L^2\pa{\R / \Z,\R^N}\to \pa{\ell^2}^N$ defined as
	$$F_N\pa{\pa{f_m}_{m=1}^N}=\pa{\left\{\widehat{f_m}(j)\right\}_{j\in\Z}}_{m=1}^N$$
	that is invertible, with inverse $F_N^{-1}:\pa{\ell^2}^N \to L^2\pa{\R / \Z,\R^N}$
	$$F_N^{-1}\pa{\pa{\left\{a^m_j\right\}_{j\in\Z}}_{m=1}^N}(x)=\pa{\sum_{j\in\Z} a^m_j e^{2\pi i j x}}_{m=1}^N.$$
	Now we compose these operators as follows
	\begin{align*}
		&\pa{F_N^{-1} \circ A_T \circ F_N}\pa{\pa{f_m}_{m=1}^N}(x)
		=\pa{F_N^{-1} \circ A_T }\pa{\pa{\left\{\widehat{f_m}(j)\right\}_{j\in\Z}}_{m=1}^N}(x) \\
		& \quad = F_N^{-1}\pa{\pa{\left\{ \sum_{s=1}^N\sum_{j\in\Z} \widehat{f_s}(j) 
		T_n\pa{\sinc}\pa{N(m-j)-s+1}\right\}_{m\in\Z}}_{n=1}^N}(x) \\
		& \quad = F_N^{-1}\pa{\pa{\left\{ \sum_{s=1}^N\sum_{j\in\Z} \widehat{f_s}(j) 
		\mathcal{F}\pa{\sum_{r\in\Z} T_n(\sinc)(Nr-s+1) e^{2\pi i r \cdot}}(m-j)\right\}_{m\in\Z}}_{n=1}^N}(x) \\
		& \quad = F_N^{-1}\pa{\pa{\left\{ \sum_{s=1}^N \mathcal{F}\pa{f_s(\cdot) 
		\sum_{r\in\Z} T_n(\sinc)(Nr-s+1) e^{2\pi i r \cdot}}(m)\right\}_{m\in\Z}}_{n=1}^N}(x) \\
		& \quad = \pa{\sum_{s=1}^N f_s(x) \pa{\sum_{r\in\Z} T_n(\sinc)(Nr-s+1) e^{2\pi i r x}}}_{n=1}^N \\
		& \quad =\pa{\sum_{r\in\Z} T_n(\sinc)(Nr-s+1) e^{2\pi i r x}}_{n,s=1}^N \pa{f_s(x)}_{s=1}^N.
	\end{align*}
	It is clear that the operator $A_T:\pa{\ell^2}^N\to\pa{\ell^2}^N$ is invertible if and only if 
	$F_N^{-1} \circ A_T \circ F_N:L^2\pa{\R / \Z,\R^N}\to L^2\pa{\R / \Z,\R^N}$ is invertible and, due 
	to what we just obtained, this is equivalent to the existence of a bounded inverse of the matrix of functions
	$$\pa{\sum_{r\in\Z} T_n(\sinc)(Nr-s+1) e^{2\pi i r x}}_{n,s=1}^N.$$
	Using Lemma \ref{lem2} we obtain that, for $x\in \pa{\frac{N-2}{2},\frac{N}{2}}$, we have that
	\begin{align*}
		&\pa{\sum_{r\in\Z} T_n(\sinc)(Nr-s+1) e^{2\pi i r x}}_{n,s=1}^N \\
		&\quad = \pa{\frac{1}{N} \sum_{m=1}^N \widehat{\sinc}\pa{\frac{m-1-x}{N}} K_n\pa{\frac{m-1-x}{N}}
		e^{-\frac{2\pi i \pa{s-1} \pa{m-1-x}}{N}}}_{n,s=1}^N \\
		&\quad = \pa{\frac{1}{N}K_n\pa{\frac{m-1-x}{N}}}_{n,m=1}^N 
        \pa{e^{-\frac{2\pi i (s-1) (m-1-x)}{N}}}_{m,s=1}^N \\
		&\quad = \pa{\mathcal{M}_T(x)}^T \pa{e^{-\frac{2\pi i (s-1) (m-1-x)}{N}}}_{m,s=1}^N 
	\end{align*}
	and as
	\begin{align*}
		\det \pa{e^{-\frac{2\pi i (s-1) (m-1-x)}{N}}}_{m,s=1}^N
		&= \prod_{1\leq n < m \leq N} \pa{e^{-\frac{2\pi i (m-1-x)}{N}}-e^{-\frac{2\pi i (n-1-x)}{N}}} \\
		&= \prod_{1\leq n < m \leq N} e^\frac{2\pi i x}{N} 
		\pa{e^{-\frac{2\pi i (m-1)}{N}}-e^{-\frac{2\pi i (n-1)}{N}}} \neq 0
	\end{align*}
	we deduce that $A_T:\pa{\ell^2}^N\to\pa{\ell^2}^N$ is invertible if and only if $\essinf_{x\in \pa{\frac{N-2}{2},\frac{N}{2}}} \abs{\det \mathcal{M}_T(x)}>0$, by Lemma \ref{lem1}. So now we will show that 2) holds if and only if $A_T:\pa{\ell^2}^N\to\pa{\ell^2}^N$ is invertible. One of the implications is easy to obtain as if 2) holds, then we have that
	$$\pa{\suc{f(Nm+n-1)}_{m\in\Z}}_{n=1}^N
	=\pa{\suc{\sum_{s=1}^N \sum_{j\in\Z} T_s(f)(Nj) g^T_s(N(m-j)+n-1)}_{m\in\Z}}_{n=1}^N$$
	and from this we can deduce that $\pa{A_T}^{-1}:\pa{\ell^2}^N\to\pa{\ell^2}^N$ exists and is defined by
	$$\pa{A_T}^{-1}\pa{\pa{\suc{b^n_m}_{m\in\Z}}_{n=1}^N}
	=\pa{\suc{\sum_{s=1}^N \sum_{j\in\Z} b_j^s \ g^T_s(N(m-j)+n-1)}_{m\in\Z}}_{n=1}^N.$$
	For the other implication, we start by assuming that $A_T:\pa{\ell^2}^N\to\pa{\ell^2}^N$ is invertible, 
	note that its inverse has to satisfy
	$$\pa{A_T}^{-1}\pa{\pa{\left\{ T_n(f)\pa{km}\right\}_{m\in\Z}}_{n=1}^N}
	=\pa{\left\{ f(km+n-1) \right\}_{m\in\Z}}_{n=1}^N,$$
	for all $f\in PW_\pi$. Now, for each $s=1,...,N$, we will consider the elements of $\pa{\ell^2}^N$ defined as
    $$\pa{\suc{C^T_{s,n,m}}_{m\in\Z}}_{n=1}^N:=\pa{A_T}^{-1}\pa{\pa{\suc{\delta_{n,s}\delta_{m,0}}_{m\in\Z}}_{n=1}^N}$$
    and use them to define the following functions in $PW_\pi$
    $$g^T_s(x)= \sum_{n=1}^N \sum_{m\in\Z} C^T_{s,n,m} \sinc (x-Nm-n+1).$$
    Using the linearity of $\pa{A_T}^{-1}$ we can obtain that
    \begin{align*}
		&\pa{\left\{ f(Nm+n-1) \right\}_{m\in\Z}}_{n=1}^N
		= \pa{A_T}^{-1}\pa{\pa{\left\{ T_n(f)(Nm)\right\}_{m\in\Z}}_{n=1}^N} \\
		&\quad = \pa{A_T}^{-1}\pa{\pa{\left\{ 
		\sum_{s=1}^N \sum_{j\in\Z} T_s(f)(Nj) \delta_{n,s}\delta_{m-j,0} \right\}_{m\in\Z}}_{n=1}^N} \\
		&\quad = \sum_{s=1}^N \sum_{j\in\Z} T_s(f)(Nj)
		\pa{A_T}^{-1}\pa{\pa{\left\{ \delta_{n,s}\delta_{m-j,0} \right\}_{m\in\Z}}_{n=1}^N} \\
		&\quad = \sum_{s=1}^N \sum_{j\in\Z} T_s(f)(Nj)
		\pa{\suc{C^T_{s,n,m-j}}_{m\in\Z}}_{n=1}^N
		= \pa{\suc{\sum_{s=1}^N \sum_{j\in\Z} T_s(f)(Nj) 
		C^T_{s,n,m-j}}_{m\in\Z}}_{n=1}^N,
	\end{align*}
	for any function $f\in PW_\pi$. Using this we can finally obtain that
	\begin{align*}
		f(x) &= \sum_{n=1}^N \sum_{m\in\Z} f(Nm+n-1) \sinc(x-Nm-n+1) \\
		&= \sum_{n=1}^N \sum_{m\in\Z} \sum_{s=1}^N \sum_{j\in\Z} T_s(f)(Nj) C^T_{s,n,m-j} \sinc(x-Nm-n+1) \\
		&= \sum_{s=1}^N \sum_{j\in\Z} T_s(f)(Nj) \sum_{n=1}^N \sum_{m\in\Z} C^T_{s,n,m-j} \sinc(x-Nm-n+1) \\
		&= \sum_{s=1}^N \sum_{j\in\Z} T_s(f)(Nj) \sum_{n=1}^N \sum_{m\in\Z} C^T_{s,n,m} \sinc(x-Nj-Nm-n+1) \\
		&= \sum_{s=1}^N \sum_{j\in\Z} T_s(f)(Nj) g^T_s(x-Nj),
	\end{align*}
	for any function $f\in PW_\pi$. So the only thing left to show is that the functions $g^T_n$, $n=1,..,N$, satisfy the characterizations i) and ii). Using the definition of these functions, we get 
	\begin{align*}
		\pa{\suc{T_n\pa{g^T_s}(Nm)}_{m\in\Z}}_{n=1}^N
		&= A_T \pa{\pa{\suc{g^T_s(Nm+n-1)}_{m\in\Z}}_{n=1}^N} \\
		&= A_T \pa{\pa{\suc{C^T_{s,n,m}}_{m\in\Z}}_{n=1}^N}
		= \pa{\suc{\delta_{n,s}\delta_{m,0}}_{m\in\Z}}_{n=1}^N,
	\end{align*}
	showing that these functions satisfy characterization i), the uniqueness comes from the existence of the interpolation formula. Using this together with Lemma \ref{lem2} we obtain
	$$ \delta_{n,s} =\sum_{m\in\Z} T_n\pa{g^T_s}(Nm) e^{2\pi i m x} \\
	= \frac{1}{N} \sum_{m=1}^N \widehat{g^T_s}\pa{\frac{m-1-x}{N}} K_n\pa{\frac{m-1-x}{N}},$$
	for almost every $x\in\pa{\frac{N-2}{2},\frac{N}{2}}$, or equivalently
	$$Id= \pa{\widehat{g^T_s}\pa{\frac{m-1-x}{N}}}_{s,m=1}^N \pa{\frac{1}{N} K_n\pa{\frac{m-1-x}{N}}}_{m,n=1}^N,$$
	for almost every $x\in\pa{\frac{N-2}{2},\frac{N}{2}}$, therefore obtaining characterization ii). \\
\end{proof}

\end{document}
